\subsection{Ausgangsspannung einer zweistufigen OPV-Schaltung bestimmen}
% Alt: Aufgabe 15
\Aufgabe{
    \label{Aufg15}
    Gegeben ist die folgende Schaltung. Berechnen Sie die Ausgangsspannung \( U_\mathrm{a} \). \\
    Die Widerstandswerte lauten: \\
    $
        R_\mathrm{1} = \mathrm{10\,k\Omega}, \
        R_\mathrm{2} = \mathrm{30\,k\Omega}, \
        R_\mathrm{3} = \mathrm{30\,k\Omega}, \
        R_\mathrm{4} = \mathrm{10\,k\Omega}
    $
    \begin{figure}[H]
        \centering
        \resizebox{0.95\textwidth}{!}{\input{Tikz/src/FigZweistufigerOPV.tex}\alttext{Das Bild zeigt eine elektronische Schaltung mit zwei Operationsverstärkern, bezeichnet als N1 und N2.
                Links befinden sich zwei Eingänge, \(U_{\mathrm{E1}}\) oben und \(U_{\mathrm{E2}}\) unten. Vom oberen Eingang \(U_{\mathrm{E1}}\) führt ein Widerstand \(R_1\) zu einem Verbindungspunkt, von dem ein Widerstand \(R_2\) zum Ausgang des Operationsverstärkers N1 zurückführt. Der Ausgang von N1 ist über einen Widerstand \(R_3\) mit dem Eingang des zweiten Operationsverstärkers N2 verbunden.
                Der untere Eingang \(U_{\mathrm{E2}}\) ist über einen weiteren Widerstand \(R_1\) mit dem gleichen Verbindungspunkt wie der obere Eingang verbunden. Von diesem Verbindungspunkt führt ein Widerstand \(R_2\) zur Masse (Erdung). Der zweite Operationsverstärker N2 besitzt ebenfalls einen Rückkopplungswiderstand \(R_4\) zwischen seinem Ausgang und seinem Eingang.
                Der Ausgang der gesamten Schaltung ist mit \(U_\mathrm{A}\) bezeichnet und befindet sich rechts außen am Ende der Schaltung. Die Masseverbindungen sind an mehreren Stellen in der Schaltung eingezeichnet.}}
        \label{fig:FigZweistufigerOPV}
    \end{figure}
}

\Loesung{
    Die Ausgangsspannung nach dem ersten Operationsverstärker (\( U_\mathrm{a1} \)) ergibt sich zu:
    \begin{align*}
        U_\mathrm{a1} & = (U_\mathrm{E1} - U_\mathrm{E2}) \cdot \frac{R_\mathrm{2}}{R_\mathrm{1}}                 \\
        U_\mathrm{a1} & = \frac{\mathrm{30\,k\Omega}}{\mathrm{10\,k\Omega}} \cdot (U_\mathrm{E1} - U_\mathrm{E2})
    \end{align*}
    Daraus ergibt sich für die gesamte Ausgangsspannung (\( U_\mathrm{a} \)):
    \begin{align*}
        U_\mathrm{a} & = - \frac{R_\mathrm{4}}{R_\mathrm{3}} \cdot U_\mathrm{a1}                                                                                           \\
        U_\mathrm{a} & = - \frac{\mathrm{10\,k\Omega}}{\mathrm{30\,k\Omega}} \cdot \frac{\mathrm{30\,k\Omega}}{\mathrm{10\,k\Omega}} \cdot (U_\mathrm{E1} - U_\mathrm{E2}) \\
        U_\mathrm{a} & = U_\mathrm{E2} - U_\mathrm{E1}
    \end{align*}


    Siehe: Abschnitt \ref{fig: Verschaltung Verstaerker} und Tabelle \ref{tab:Grundschaltungen}
}